%***************************************************************************************************
% preamble/headings.tex - formatowanie tytułów rozdziałów i sekcji, głębokość numeracji i spisu treści.
%***************************************************************************************************

\setcounter{secnumdepth}{3}    % numerujemy do poziomu \subsubsection
\setcounter{tocdepth}{2}       % w spisie treści do poziomu \subsection

% Etykieta przed tytułem rozdziału (w załącznikach zmieniana przez \StartAppendices)
\newcommand{\ChapterLabel}{\thechapter.}

\titleformat{\chapter}[hang]
  {\normalfont\huge\bfseries}{\ChapterLabel}{1ex}{}
\titlespacing*{\chapter}{0pt}{0pt}{25pt}

\titleformat{\section}[hang]
  {\normalfont\Large\bfseries}{\thesection.}{1ex}{}
\titlespacing*{\section}{0pt}{3.5ex plus 1ex minus .2ex}{2ex plus .2ex}

\titleformat{\subsection}[hang]
  {\normalfont\large\bfseries}{\thesubsection.}{1ex}{}
\titlespacing*{\subsection}{0pt}{3ex plus 1ex minus .2ex}{1.5ex plus .2ex}

\titleformat{\subsubsection}[hang]
  {\normalfont\normalsize\bfseries}{\thesubsubsection.}{1ex}{}
\titlespacing*{\subsubsection}{0pt}{2.5ex plus 1ex minus .2ex}{1ex plus .2ex}

% Brak wdów i sierot (pojedynczych wierszy akapitu na początku/końcu strony)
\clubpenalty=10000
\widowpenalty=10000

% Zalecana polska typografia: brak "samotnych" spójników na końcu wiersza ustawia się ręcznie
% za pomocą twardej spacji (np. i~inne, w~tabeli).

% Załączniki: \StartAppendices przełącza numerację na literową i etykietę "Załącznik A."
\newcommand{\StartAppendices}{%
  \clearpage
  \appendix
  \renewcommand{\ChapterLabel}{\appendixname~\thechapter.}%
}
